% -----------------------------------------------------------------------------
% Strict-SI reduced momentum-space rules used by all scattering examples
% -----------------------------------------------------------------------------

进入具体散射例子以前，把由 SI 作用量、SI 二点函数和 LSZ 外腿归一化共同产生的常数统一提出，定义一套\emph{约化图因子}。约化后的内部线与顶角仍以物理四动量表示
\begin{equation}
 p^\mu=(E_{\mathbf p}/c,\mathbf p),
 \qquad p^2=m^2c^2,
 \qquad [p^\mu]=\mathrm{kg\,m\,s^{-1}},
 \label{eq:reduced-si-physical-p-r10}
\end{equation}
质量从不写成 ``$m$ 具有动量量纲''；凡是与 $p^\mu$ 相加的质量项一律写成 $mc$。

对电荷 $q_a$ 定义
\begin{equation}
 g_a\equiv\frac{q_a}{\sqrt{\varepsilon_0\hbar c}},
 \qquad
 [g_a]=1.
 \label{eq:si-reduced-charge-coupling-r10}
\end{equation}
特别地，电子满足
\begin{equation}
 q_{\mathrm e}=-e,
 \qquad
 g_e=-g_{\rm em},
 \qquad
 g_{\rm em}\equiv\frac{e}{\sqrt{\varepsilon_0\hbar c}}
 =\sqrt{4\pi\alpha},
 \qquad
 \alpha=\frac{e^2}{4\pi\varepsilon_0\hbar c}.
 \label{eq:gem-alpha-si-r10}
\end{equation}
因此 $e$ 始终是以 C 为单位的基本电荷，而 $g_{\rm em}$ 才是图计算中出现的无量纲 QED 耦合。后续计算绝不再用同一个符号 $e$ 同时表示 ``库仑'' 和 ``无量纲顶角常数''。


对 Yukawa 理论，前面使用的 SI 标量场归一化给 $[\varphi]=\sqrt{\mathrm{J/m}}$、$[\psi]=\mathrm{m^{-3/2}}$。于是 $\varphi\bar\psi\psi$ 前的系数必须具有 $\sqrt{\mathrm{J\,m}}=\sqrt{\hbar c}$ 的量纲。把无量纲 Yukawa 耦合统一记作 $y$，则相互作用写成
\begin{equation}
 \boxed{
 \Lag_Y^{\rm int}=-y\sqrt{\hbar c}\,\varphi\bar\psi\psi,}
 \qquad [y]=1.
 \label{eq:yukawa-si-coupling-r10}
\end{equation}
若有不同费米子种类，则分别使用 $y_A,y_B$。对于实标量四次项，$[\varphi^4]=\mathrm{J^2m^{-2}}$，要恢复拉氏密度单位 $\mathrm{Jm^{-3}}$，系数必须带 $1/(\hbar c)$。因此把剩余无量纲系数定义为 $g_4$，相互作用写成
\begin{equation}
 \boxed{
 \Lag_{\varphi^4}^{\rm int}
 =-\frac{g_4}{4!\,\hbar c}\,\varphi^4,}
 \qquad [g_4]=1,
 \label{eq:phi4-si-coupling-r10}
\end{equation}
因为 $[\varphi^4]=\mathrm{J^2/m^2}$，乘 $1/(\hbar c)$ 后正好恢复能量密度量纲 $\mathrm{J/m^3}$。

前面的物理二点函数都带有由场归一化固定的 $\hbar,c,\mu_0$ 因子。为把这些公共 SI 归一化因子从散射图的拓扑计算中分离出来，定义一组与无量纲截肢振幅 $\mathcal M$ 配套的\emph{约化内部线}。它们只是计算核的记号定义；与约化顶角、LSZ 外腿和相空间重新组合后，必须恢复原来的 SI $S$ 矩阵与可观测截面。完整一光子交换、Yukawa 交换和 $\varphi^4$ 例子将直接检验这一组合。三类约化内部线定义为：
\begin{align}
 \Delta_{\rm red}(p)
 &\equiv\frac{\ii}{p^2-m^2c^2+\ii0^+},\\
 S_{\rm red}(p)
 &\equiv\frac{\ii(\slashed p+mc)}
 {p^2-m^2c^2+\ii0^+},
 \\
 D_{{\rm red},F}^{\mu\nu}(q)
 &\equiv\frac{-\ii\eta^{\mu\nu}}
 {q^2+\ii0^+}.
 \label{eq:si-reduced-propagators-r10}
\end{align}
标量与光子约化传播子的量纲均为动量的负二次方，Dirac 约化传播子由于分子含一个动量因子而具有动量的负一次方。三者的极点位置和 Feynman $+\ii0^+$ 边界条件与物理二点函数完全相同。

同样地，下列“约化顶角”也是与这套 $\mathcal M$ 定义配套的计算核，而不是另一套相互作用 Hamiltonian。它们把物理作用量中的 SI 常数与内部线、外腿归一化成组提出后，只留下无量纲耦合和真正决定自旋/内部指标结构的矩阵：
\begin{equation}
 \boxed{
 \varphi^4:\ -\ii g_4,
 \qquad
 \varphi\bar\psi\psi:\ -\ii y,
 \qquad
 A_\mu\bar\psi\psi:\ -\ii g_a\gamma^\mu.}
 \label{eq:si-reduced-vertices-r10}
\end{equation}
约化以后得到的截肢 $2\to2$ 振幅 $\mathcal M$ 无量纲，截面中的面积量纲由 Lorentz 不变相空间与入射通量组合中的 $\hbar^2/s$ 提供。

Wick 收缩还固定三条图计算的组合学规则。对四点 $\varphi^4$ 泡图，两条相同内部线的交换不产生新图，因此该拓扑带
\begin{equation}
 S_{\rm bubble}=\frac12.
 \label{eq:phi4-bubble-symmetry-factor-dp68}
\end{equation}
对两个全同外部费米子，交换末态标签需要一次费米算符奇置换，所以相干振幅按
\begin{equation}
 \mathcal M_{ff}=\mathcal M_t-\mathcal M_u
 \label{eq:identical-fermion-exchange-sign-dp68}
\end{equation}
组合。若 Wick 收缩形成一条闭合费米环，同一个 Grassmann 奇偶性给出
\begin{equation}
 s_{\rm loop}=-1,\qquad \text{closed fermion loop: }-\operatorname{Tr}(\cdots),
 \label{eq:closed-fermion-loop-sign-dp68}
\end{equation}
即每条独立闭合费米环额外贡献一个负号；这正是通用费米 Wick 结果~\eqnrefs{eq:generic-closed-fermion-loop-r68} 在 QED 收缩中的专门化。泡图的 $1/2$ 与全同外腿交换负号都直接来自同一 Dyson--Wick 收缩计数。

在圈图中还需要固定一个不会偷偷改变单位的积分记号。若 $\ell^\mu$ 是物理四动量，定义
\begin{equation}
 \boxed{
 \int_{\ell}^{(d)}F(\ell)
 \equiv
 \mu_{\rm DR}^{4-d}
 \int\frac{\dd^d \ell}{(2\pi)^d}F(\ell),}
 \label{eq:reduced-loop-measure-si-r10}
\end{equation}
其中 $\mu_{\rm DR}$ 本身具有\emph{动量}量纲。这里的 $(2\pi)^{-d}$ 来自把物理 Fourier 测度 $(2\pi\hbar)^{-d}\dd^d \ell$ 中与约化传播子/顶角归一化成对出现的 $\hbar^d$ 一并提出；所以 $\ell$ 始终以 $\mathrm{kg\,m\,s^{-1}}$ 为单位。在 $d=4$ 时 $\int_\ell^{(4)}$ 具有动量四次方量纲。
